\documentclass{beamer}

\usepackage[russian]{babel}
\usepackage[utf8]{inputenc}
\usepackage{cmap}
\usepackage{graphicx}
\usepackage{xspace}
\usepackage{psfrag}
\usepackage{hyperref}
\usepackage[normalem]{ulem}
\usepackage{listings}
\usepackage{colortbl}

\beamertemplatenavigationsymbolsempty
\setbeamertemplate{footline}[frame number]
\usecolortheme{seahorse}

%%Исправляем греческие буквы в формулах, они должны быть прямыми, а не курсивными
\DeclareSymbolFont{greek}{U}{eur}{m}{n}                                                        
\SetSymbolFont{greek}{bold}{U}{eur}{b}{n}                                                      
\DeclareSymbolFontAlphabet{\gr}{greek}                                                         
\DeclareMathSymbol{\alpha}{\mathord}{greek}{"0B}                                               
\DeclareMathSymbol{\beta}{\mathord}{greek}{"0C}                                                
\DeclareMathSymbol{\gamma}{\mathord}{greek}{"0D}                                               
\DeclareMathSymbol{\delta}{\mathord}{greek}{"0E}                                               
\DeclareMathSymbol{\epsilon}{\mathord}{greek}{"0F}                                             
\DeclareMathSymbol{\zeta}{\mathord}{greek}{"10}                                                
\DeclareMathSymbol{\eta}{\mathord}{greek}{"11}                                                 
\DeclareMathSymbol{\theta}{\mathord}{greek}{"12}                                               
\DeclareMathSymbol{\iota}{\mathord}{greek}{"13}                                                
\DeclareMathSymbol{\kappa}{\mathord}{greek}{"14}                                               
\DeclareMathSymbol{\lambda}{\mathord}{greek}{"15}                                              
\DeclareMathSymbol{\mu}{\mathord}{greek}{"16}                                                  
\DeclareMathSymbol{\nu}{\mathord}{greek}{"17}                                                  
\DeclareMathSymbol{\xi}{\mathord}{greek}{"18}                                                  
\DeclareMathSymbol{\pi}{\mathord}{greek}{"19}                                                  
\DeclareMathSymbol{\rho}{\mathord}{greek}{"1A}                                                 
\DeclareMathSymbol{\sigma}{\mathord}{greek}{"1B}                                               
\DeclareMathSymbol{\tau}{\mathord}{greek}{"1C}                                                 
\DeclareMathSymbol{\upsilon}{\mathord}{greek}{"1D}                                             
\DeclareMathSymbol{\phi}{\mathord}{greek}{"1E}                                                 
\DeclareMathSymbol{\chi}{\mathord}{greek}{"1F}                                                 
\DeclareMathSymbol{\psi}{\mathord}{greek}{"20}                                                 
\DeclareMathSymbol{\omega}{\mathord}{greek}{"21}                                               
\DeclareMathSymbol{\varepsilon}{\mathord}{greek}{"22}                                          
\DeclareMathSymbol{\vartheta}{\mathord}{greek}{"23}                                            
\DeclareMathSymbol{\varpi}{\mathord}{greek}{"24}                                               
\let\varrho\rho                                                                                
\let\varsigma\sigma                                                                            
\DeclareMathSymbol{\varphi}{\mathord}{greek}{"27}

\title{Развёртывание и исследование ARM64 Linux-системы на x86\_64-платформе с использованием симулятора gem5 для решения задач кросскомпиляции}
\author{Студент: Гельдыев~Б.~932204 \\ Науч.рук.: Пожидаев~Михаил~Сергеевич}
\date{}

\begin{document}
\maketitle

% Slide 2
\begin{frame}
\frametitle{Цели и задачи}
\textbf{Цель работы:} разработать и развернуть воспроизводимую среду выполнения ARM64
Linux-системы на x86\_64-хосте с использованием архитектурного симулятора gem5,
обеспечивающую интерактивную работу и возможность дальнейших экспериментов.

\medskip
\textbf{Задачи работы:}
\begin{enumerate}
    \item Проанализировать существующие подходы к запуску ARM-систем на x86\_64
    \item Изучить архитектуру и возможности симулятора gem5
    \item Подготовить окружение для full-system симуляции ARM64
    \item Загрузить ARM64 Linux kernel и root filesystem (Ubuntu)
    \item Обеспечить интерактивный доступ к системе через консоль
    \item Проверить корректность загрузки и работы пользовательского пространства
    \item Подготовить основу для автоматизации и контейнеризации решения
\end{enumerate}
\end{frame}

% Slide 3
\begin{frame}
\frametitle{Функциональные требования}
\begin{enumerate}
    \item Запуск ARM64 Linux Kernel в режиме full-system симуляции
    \item Поддержку загрузки с виртуального диска
    \item Корректное монтирование корневой файловой системы
    \item Поддержка стандартных Linux-интерфейсов:
    \begin{itemize}
        \item Файловая система
        \item Процессы
        \item Консольный ввод/вывод
    \end{itemize}
    \item Возможность интерактивной работы пользователя (shell)
    \item Поддержку автоматического выполнения сценариев при загрузке
    \item Возможность изменения конфигурации симулируемой системы:
    \begin{itemize}
        \item Тип CPU
        \item Объём памяти
        \item Количество ядер
    \end{itemize}
\end{enumerate}
\end{frame}

% Slide 4
\begin{frame}
\frametitle{Нефункциональные требования}
\begin{enumerate}
    \item Повторный запуск должен приводить к одинаковым результатам при одинаковой конфигурации
    \item Возможность запуска на стандартной системе x86\_64 Linux-системе без специализированного оборудования
    \item Возможность дальнейшего добавления:
    \begin{itemize}
        \item Архитектурных экспериментов
        \item Новых образов ОС
        \item Других моделей CPU
    \end{itemize}
    \item Возможность получения отладочной информации и логов работы системы
    \item Поддержка и запуск без ручного вмешательства (контейнеры)
\end{enumerate}
\end{frame}

% Slide 5
\begin{frame}
\frametitle{Обзор существующих подходов}
\begin{table}
\small
\begin{tabular}{|p{2.5cm}|p{2cm}|p{1.8cm}|p{2.5cm}|}
\hline
\rowcolor{blue!20} Критерий & KVM / VirtualBox & QEMU & gem5 \\
\hline
Запуск ARM на x86\_64 & Нет & Да & Да \\
\hline
Тип технологии & Виртуализация & Эмуляция & Архитектурная симуляция \\
\hline
Режим full-system & Нет & Да & Да \\
\hline
Точность моделирования CPU & Низкая & Средняя & Высокая \\
\hline
Моделирование микроархитектуры & Нет & Ограничено & Да \\
\hline
Поддержка архитектурных экспериментов & Нет & Ограничено & Да \\
\hline
\end{tabular}
\end{table}
\end{frame}

% Slide 6
\begin{frame}
\frametitle{Обоснование выбора симулятора gem5}
\begin{enumerate}
    \item gem5 позволяет запускать ARM64 Linux-системы на x86\_64-хосте в режиме full-system симуляции
    \item В отличие от QEMU, gem5 моделирует внутреннюю архитектуру процессора, включая:
    \begin{itemize}
        \item Конвейер
        \item Кэш-память
        \item Подсистему памяти
        \item Взаимодействие с устройствами
    \end{itemize}
    \item Возможность изменять параметры моделируемой системы:
    \begin{itemize}
        \item Тип процессорной модели
        \item Количество ядер
        \item Объём и тип памяти
    \end{itemize}
\end{enumerate}
\medskip
\textbf{Ограничение:} низкая производительность по сравнению с эмуляцией и виртуализацией
\end{frame}

% Slide 7
\begin{frame}
\frametitle{Общая архитектура}
Разрабатываемое решение представляет собой многоуровневую систему, в которой ARM64-платформа полностью моделируется поверх x86\_64-хоста.
\begin{enumerate}
    \item \textbf{Хост система:} архитектура x86\_64, ОС Linux; запуск эмулятора, управление конфигурацией, подключение к консоли гостевой системы
    \item \textbf{Архитектурный симулятор gem5:} полная симуляция аппаратной платформы: ARM64 CPU, ОЗУ, устройства ввода-вывода, виртуальные шины
    \item \textbf{Симулируемая ARM64-система:} процессор ARMv8-A, ядро vmlinux, корневая ФС Ubuntu (ext4), виртуальный диск (virtio), консоль PL011 UART
    \item \textbf{Пользовательский интерфейс:} подключение через ttyAMA0, TCP-порты, проброшенные gem5, поддержка интерактивной shell-сессии
\end{enumerate}
\end{frame}

% Slide 8
\begin{frame}[fragile]
\frametitle{Подготовка окружения}
\begin{lstlisting}[basicstyle=\small\ttfamily,frame=none]
$M5_PATH/
|-- binaries/
|   |-- boot.arm
|   |-- boot.arm64
|   `-- vmlinux.arm64
`-- disks/
    `-- ubuntu-18.04-arm64-docker.img
\end{lstlisting}
\medskip
\textbf{Переменная окружения:}
\begin{lstlisting}[basicstyle=\small\ttfamily,frame=none]
EXPORT M5_PATH=$HOME/m5/system
\end{lstlisting}
\end{frame}

% Slide 9
\begin{frame}[fragile]
\frametitle{Запуск симуляции ARM64 Linux}
\begin{columns}
\begin{column}{0.57\textwidth}
\begin{lstlisting}[basicstyle=\tiny\ttfamily,frame=single]
./build/ARM/gem5.opt \
  configs/example/arm/starter_fs.py \
  --kernel \
  "$M5_PATH/binaries/vmlinux.arm64" \
  --disk-image \
  "$M5_PATH/disks/ubuntu-18.04-arm64-docker.img" \
  --cpu atomic \
  --num-cores 1 \
  --mem-size 2GB
\end{lstlisting}
\end{column}
\begin{column}{0.43\textwidth}
Параметры запуска:

\medskip
\texttt{--cpu atomic} --- быстрый старт и отладка

\medskip
\texttt{--num-cores 1} --- однопроцессорная конфигурация

\medskip
\texttt{--mem-size 2GB} --- ограничение объёма памяти

\medskip
\texttt{starter\_fs.py} --- режим full-system
\end{column}
\end{columns}
\end{frame}

% Slide 10
\begin{frame}[fragile]
\frametitle{Подключение и результат работы}
\begin{lstlisting}[basicstyle=\small\ttfamily,frame=none]
rlwrap -cAr nc 127.0.0.1 3456
\end{lstlisting}
\begin{lstlisting}[basicstyle=\tiny\ttfamily,
  backgroundcolor=\color{black},
  basicstyle=\tiny\ttfamily\color{white},
  frame=single,framesep=4pt]
Ubuntu 18.04.2 LTS aarch64-gem5 ttyAMA0

aarch64-gem5 login: root (automatic login)

Welcome to Ubuntu 18.04.2 LTS (GNU/Linux 4.18.0+ aarch64)

 * Documentation:  https://help.ubuntu.com
 * Management:     https://landscape.canonical.com
 * Support:        https://ubuntu.com/advantage

root@aarch64-gem5:~#
\end{lstlisting}
\end{frame}

% Slide 11
\begin{frame}[shrink=5]
\frametitle{Процесс загрузки ARM64 Linux в gem5}
\begin{enumerate}
    \item \textbf{Запуск симулятора gem5} --- инициализация симулируемой аппаратной платформы. Создание виртуального ARM64-окружения (CPU, память, устройства)
    \item \textbf{Загрузка bootloader} --- используются boot.arm и boot.arm64. Подготовка процессора к передаче управления ядру Linux
    \item \textbf{Запуск ядра Linux} --- ARM64 ядро (vmlinux.arm64) загружается в память. Передача параметров командной строки ядра. Инициализация драйверов и подсистем ядра
    \item \textbf{Инициализация устройств} --- последовательная консоль (PL011, ttyAMA0), виртуальный диск (virtio-blk), таймеры, прерывания, память
    \item \textbf{Монтирование корневой файловой системы} --- обнаружение /dev/vda1, монтирование rootfs (ext4), переход в пользовательское пространство
    \item \textbf{Запуск userland} --- выполнение init-скриптов, запуск getty на ttyAMA0, автоматический вход пользователя root
\end{enumerate}
\end{frame}

% Slide 12
\begin{frame}
\frametitle{Выводы}
\begin{enumerate}
    \item Развёрнута полноценная ARM64 Linux-система на x86\_64-хосте
    \item Использован режим full-system симуляции
    \item Обеспечена интерактивная работа через консоль
    \item Подтверждена корректная загрузка ядра и rootfs Ubuntu
\end{enumerate}
\medskip
\textbf{Направления дальнейшего развития}
\begin{enumerate}
    \item Использование более детальных моделей CPU (O3)
    \item Проведение архитектурных экспериментов с кэш-памятью
    \item Интеграция с Docker для автоматизации запусков
\end{enumerate}
\end{frame}

\end{document}
